\section{等周不等式与单调性公式}
\label{sec:isoperimetric-monotonicity}
回忆推论 7.2.6：对殆凯勒三元组 $(M,\omega,J)$，“\emph{任意
$J$-全纯曲线的像关于相伴度量 $g=\omega(\cdot,J\cdot)$ 都是极小曲面}”。我们先把引号中的陈述说得更精确。为此必须研究 $J$-全纯曲线之像的结构。下述命题是第 8.8 节相似性原理（定理 8.8.2）的直接推论，尤其参见推论 8.8.6 与定理 8.8.9；相关结果也见 (Mc87)、(Sik94)。
\begin{proposition}\label{prop:ch8-isolated-coincidences-critical-points}
设 $(N,J)$ 为任意赫尔德类 $C^\epsilon$（$\epsilon>0$）的殆复流形，并令
$\Sigma$ 为闭黎曼曲面或带边黎曼曲面。
\begin{enumerate}[label=(\arabic*)]
\item 对任何两个不同的 $J$-全纯映射 $f,g:\Sigma\to N$，满足
$f(z)=g(z)$ 的点集没有聚点。
\item 再假设 $J$ 属于 $C^1$，并令 $f:\Sigma\to N$ 为非恒定
$J$-全纯映射，则 $f$ 的临界点集合是离散的。
\end{enumerate}
同样的陈述也适用于带非空边界的曲面以及满足全实边界条件
$f:(\Sigma,\partial\Sigma)\to(N,R)$ 的 $J$-全纯映射。
\end{proposition}
此命题的一个立即推论是：当 $\Sigma$ 闭时，任意 $J$-全纯映射的像只有有限个奇点。
\begin{definition}\label{def:ch8-somewhere-injective}
若存在 $z\in\Sigma$ 使 $\#f^{-1}(f(z))=1$ 且 $df(z)\ne0$，则称
$J$-全纯曲线 $f:\Sigma\to N$ \emph{在某处单射}。
\end{definition}
下面的结构定理是伪全纯曲线研究中的重要工具；其证明见 (Mc87)。
\begin{theorem}[结构定理 (Mc87)]\label{thm:ch8-structure-somewhere-injective}
若闭黎曼曲面上的 $J$-全纯映射 $f$ 非在某处单射，则存在在某处单射的
$\bar f:\bar\Sigma\to N$ 及分歧全纯覆叠
$\phi:\Sigma\to\bar\Sigma$，使 $f=\bar f\circ\phi$。
\end{theorem}
现在可以描述在某处单射曲线的像之结构。
\begin{definition}\label{def:ch8-singular-image-point}
设 $\Sigma$ 为闭黎曼曲面，$f:\Sigma\to N$ 为在某处单射的
$J$-全纯映射，并记 $A=\operatorname{Im}f$。若 $\#f^{-1}(x)\geq2$，或
$x=f(z)$ 且 $df(z)=0$，称 $x\in A$ 为\emph{奇异}点；否则称\emph{正则}点。
\end{definition}
\begin{corollary}\label{cor:ch8-area-of-image-well-defined}
若 $\Sigma$ 是边界为 $\partial\Sigma$ 的紧黎曼曲面，且
$f:\Sigma\to N$ 为 $J$-全纯映射，记 $A=\operatorname{Im}f$。若 $A$ 有正则点，则
$\#\operatorname{Sing}(A)<\infty$，并且对 $N$ 上任意给定的度量 $g$，
$A$ 的面积良定义。
\end{corollary}
\begin{proof}
由 $\Sigma$ 的紧性与命题 8.7.1，立即有
$\#\operatorname{Sing}(A)<\infty$。每个正则点附近，$A$ 是可微嵌入曲面；故定义
\[
 \operatorname{Area}A=\operatorname{Area}(A\setminus\operatorname{Sing}(A))
 =\int_\Sigma|f_x\wedge f_y|\,dx\,dy,
\]
其中 $(x,y)$ 是 $(\Sigma,j)$ 的坐标，且
$|f_x\wedge f_y|^2=|f_x|^2|f_y|^2-\langle f_x,f_y\rangle^2$。由
$\Sigma$ 紧及 $f$ 可微，积分可积且定义良好。
\end{proof}
当 $g=g_J=\omega(\cdot,J\cdot)$ 时有 $E_J(f)=\operatorname{Area}A$。有了这些准备，我们可以陈述 $J$-全纯曲线最重要的两个几何性质：\emph{等周不等式}与\emph{单调性}。二者在可去奇点定理以及 $J$-全纯曲线模空间紧化中都扮演关键角色。先从欧氏空间中极小曲面的著名等周不等式与单调性公式开始；经典例子见 (Law80)。
\begin{lemma}[经典等周不等式]\label{lem:ch8-classical-isoperimetric}
\stfhIndex{I172}{deng^zhou^bu^deng^shi}{等周不等式（isoperimetric inequality）}
\stfhSubIndex{I173}{deng^zhou^bu^deng^shi}{等周不等式（isoperimetric inequality）}{jing^dian^de}{经典的（classical）}
\stfhIndex{I203}{mai^ke^da^fu}{麦克达夫（McDuff）}
若 $\Sigma'$ 解以嵌入可求长若尔当曲线 $\ell\subset\mathbb R^k$ 为边界的普拉托问题，则存在与 $\Sigma'$ 无关的常数 $C_1>0$，使
\begin{equation}\operatorname{Area}(\Sigma')\leq C_1^2\operatorname{length}^2(\ell).\label{eq:ch8-classical-isoperimetric}\end{equation}
\end{lemma}

为把这一经典结果转化为 $J$-全纯曲线的结果，我们采用西科拉夫 (Sik94) 所用的灵活几何设置。
\stfhIndex{I246}{xi^ke^la^fu}{西科拉夫（Sikorav）}
\begin{definition}\label{def:ch8-tame-almost-complex-triple}
若 $g$ 完备、曲率有界、单射半径有正下界，且 $J$ 关于 $g$ 一致连续，则称 $(N,J,g)$ \emph{驯服}。
\end{definition}
这里对 $J$ 只需很低的正则性；例如任意 $C^\epsilon$（$\epsilon>0$）赫尔德类已经足够。下述引理由 (Sik94) 借来。
\begin{lemma}\label{lem:ch8-tame-local-geometry}
若 $(N,J,g)$ 驯服，则：
\begin{enumerate}[label=(T\arabic*)]
\item 对所有 $x\in N$，指数映射
$\exp_x:B(0,r_0)\subset T_xN\to B(x,r_0)\subset N$ 为微分同胚；
\item 每条包含于 $B=B(x,r)$、$r\leq r_0$ 的环路 $\gamma$，都在
$B$ 内界住面积小于 $C_3\operatorname{length}(\gamma)^2$ 的圆盘；
\item 对每个 $B=B(x,r_0)$，存在辛形式 $\omega_x$，使
$\|\omega_x\|\leq1$ 且
$|X|^2\leq C_4^2\omega_x(X,JX)$。后一性质称为驯服性质，而
$\omega_x$ 称为 $(N,J)$ 的（局部）驯服形式。
\stfhIndex{I283}{xun^fu^xing^shi}{驯服形式（taming form）}
\end{enumerate}
\end{lemma}
\begin{proof}
(T1) 立即来自单射半径下界。证明 (T2)：由 $(N,J,g)$ 的假设，单射半径为正。
固定小于单射半径的常数 $r_0$。于是可取 $g$ 的一个开测地球
$B(x,r_0)$ 及一个驯服形式 $\eta_0$，使 $f(\Sigma)\subset B(x,r_0)$，并定义
\begin{equation}g_x=\eta_0(\cdot,J_x\cdot)\label{eq:ch8-local-taming-metric}\end{equation}
令 $f:D\to N$ 是像落在该球内的 $J$-全纯圆盘。由引理 8.7.6，存在紧曲面
$S$，其边界为 $f(\partial\Sigma)$，并且
\begin{equation}\operatorname{Area}_{g_0}(S)\leq C_1^2\operatorname{length}_{g_0}^2(\partial S)\label{eq:ch8-local-filling-area}\end{equation}
这里 $g_0=(\exp_x^{-1})^*g$。另一方面，
\[
\operatorname{Area}_{g_0}(S)=\int_{\Sigma'}\eta_0.
\]
由 $B$ 可缩，$[f(\Sigma)\#\overline\Sigma]=0\in H_2(N)$；结合 $f$ 的重数假设，
\[
 \int_S\eta_0=\int_Df^*\eta_0=\operatorname{Area}_{g_0}(f).
\]
度量一致比较为 $g_0/C_2^2\leq g\leq C_2^2g_0$，故
$\operatorname{Area}_{g_0}(f)\geq C_2^{-2}\operatorname{Area}_g(f)$ 且
$\operatorname{length}_{g_0}^2(\partial S)\leq C_2^{-2}
\operatorname{length}_g^2(\partial S)$。代入上式得到 (T2)，其中
$C_3^2=C_1^2/C_2^2$。

最后证明 (T3)。在 $T_xN$ 上取满足 $\eta_0(\cdot,J_x\cdot)$ 正定的常系数
$2$-形式 $\eta_0$。由 $g$ 的假设以及 $J$ 关于 $g$ 一致连续，
$\eta=(\exp_x^{-1})^*\eta_0$ 在与 $x$ 无关的小球 $B(x,r_0)$ 上良定义、为辛形式且驯服 $J$。这证明 (T3)。
\end{proof}
现在假设定义 8.7.7 中的 (T1)--(T3)，并证明一般等周不等式。在
$J$-全纯曲线研究中使用驯服形式的想法最初归功于格罗莫夫 (Gr85)。
\begin{proposition}[等周不等式]\label{prop:ch8-j-holomorphic-isoperimetric}
\stfhIndex{I172}{deng^zhou^bu^deng^shi}{等周不等式（isoperimetric inequality）}
\stfhIndex{I206}{dan^diao^xing^gong^shi}{单调性公式（monotonicity formula）}
设 $(N,J,g')$ 是满足 (T1)--(T3) 的殆埃尔米特流形，且
$r_0,C_3,C_4$ 如引理 8.7.8。若 $f:\Sigma\to N$ 是紧黎曼曲面上的
$J$-全纯映射，且在边界 $\partial\Sigma$ 附近为嵌入；又若
$f(\Sigma)\subset B(x,r_0)$，则存在只依赖于 $(N,J,g')$ 的 $C_5>0$ 使
\begin{equation}\operatorname{Area}_{g'}(f(\Sigma))\leq C_5^2\operatorname{length}_{g'}^2(\partial f(\Sigma)).\label{eq:ch8-j-holomorphic-isoperimetric}\end{equation}
\end{proposition}
\begin{proof}
由 (T2)，$f(\partial\Sigma)$ 在 $B(x,r_0)$ 内界住曲面 $S$，且其面积不超过
$C_3^2\operatorname{length}_{g'}^2(\partial f(\Sigma))$。因 $B$ 可缩，
\[
 \int_\Sigma f^*\omega_x=\int_S\omega_x
 \leq C_3^2\operatorname{length}_{g'}^2(\partial f(\Sigma)),
\]
这里用了 $\|\omega_x\|\leq1$。另一方面，(T3) 的驯服性质以及 $f$ 为
$J$-全纯给出
\[
 \operatorname{Area}_{g'}(f(\Sigma))
 \leq C_4^2\int_\Sigma f^*\omega_x.
\]
合并两式，取 $C_5=C_3C_4$ 即得。
\end{proof}
下面证明单调性；粗略地说，它断言：“若 $J$-全纯映射 $f$ 的像之直径大，
那么它的面积必定也大。”
\begin{proposition}[单调性]\label{prop:ch8-monotonicity}
设 $(N,J,g')$ 与 $r_0$ 如命题 8.7.9。令 $f:\Sigma\to N$ 是紧带边
黎曼曲面上的 $J$-全纯映射，满足
$f(\Sigma)\subset B=B(x,r)$、$f(\partial\Sigma)\subset\partial B$ 且
$x\in f(\Sigma)$。则存在使用上述同一个 $r_0$ 的 $C_7>0$，使
\begin{equation}\operatorname{Area}_{g'}(f(\Sigma))\geq C_6^2r^2.\label{eq:ch8-monotonicity}\end{equation}
\end{proposition}
\begin{proof}
令 $\Sigma_t=f^{-1}(B(x,t))$、$A(t)=\operatorname{Area}(f|_{\Sigma_t})$。因
$f\in C^{r+1}$，对几乎所有正则值 $t>0$，萨德定理说明 $\Sigma_t$ 是紧子曲面（不必连通，但只有有限个连通分支），且 $C^1$ 边界为
$\partial\Sigma_t=f^{-1}(\partial B(x,t))$。记
$\ell(t)=\operatorname{length}(f(\partial\Sigma_t))$。于是 $A(t)$ 绝对连续且
$A'(t)=\ell(t)$ 几乎处处。对每个 $\Sigma_t$ 应用等周不等式，
$A(t)\leq C_5^2\ell(t)^2$，从而 $A'(t)\geq\sqrt{A(t)}/C_5$。这等价于
\[
 (\sqrt{A(t)})'\geq\frac1{2C_5}
\]
在 $A$ 可微时成立。从 $0$ 到 $r$ 积分，
$\sqrt{A(r)}\geq r/(2C_5)$，即
$A(r)\geq r^2/(4C_5^2)$。取 $C_6=1/(2C_5)$ 即得。（可将此证明与定理 8.6.1 比较。）
\end{proof}
